% REJECTED BEFORE ADMISSION. Preserved first draft; DO NOT include in the paper.
% Error: nominal posterior-minimum score was treated as actual continuing risk.
\section{Separating Identification Value, State Value, and the Paid Horizon}
The exact initialization-minimax envelope can be compared with the complete
full-budget fixed catalogue without confusing three distinct effects.
For one fixed probe, let $I_{\rm full}$ be the best full-budget fixed risk,
$I_{\rm post}=\min_a\sum_j\pi_jJ_{j,a}^{r}$ the best observation-independent
program after paying the probe, $R_{\rm marg}=\mathbb E_g\min_aF_a(Y)$,
and $R_{\rm rob}=\mathbb E_g\min_aL_a(Y)$.
All post-probe quantities use exactly the same remaining budgets and catalogue,
including stop.
Define
\[
C=I_{\rm post}-I_{\rm full},\quad
V_{\rm reg}=I_{\rm post}-R_{\rm marg},\quad
V_{\rm state}=R_{\rm marg}-R_{\rm rob}.
\]
Then
\[
I_{\rm full}-R_{\rm rob}=V_{\rm reg}+V_{\rm state}-C.
\]
Here $V_{\rm reg},V_{\rm state}\ge0$ by inclusion of the corresponding selector
classes. The quantity $C$ is the exact, possibly signed, change in the best
fixed risk caused by the paid remaining horizon; no monotonicity of the affine
learner in its horizon is assumed.
This decomposition is an accounting identity using classical decision value,
not a new general information-value theorem.
Let $K$ be the finite state-value ceiling proved above, and define
$I_{\rm oracle}=\sum_j\pi_j\min_aJ_{j,a}^{r}$.
Because a posterior-weighted average is at least the posterior average of the
individual regime minima, $R_{\rm marg}\ge I_{\rm oracle}$.
Consequently
\[
0\le V_{\rm state}\le K,\quad
0\le V_{\rm reg}\le I_{\rm post}-I_{\rm oracle},\quad
I_{\rm full}-R_{\rm rob}\le I_{\rm full}-I_{\rm oracle}+K.
\]
Thus $I_{\rm oracle}-K\ge I_{\rm full}$ rules out any strict improvement by
any randomized measurable selector following this fixed probe, even if its
regime inference is optimal.
For a desired relative gain $\epsilon>0$, the stronger necessary condition is
$I_{\rm full}-I_{\rm oracle}+K\ge\epsilon I_{\rm full}$.
Conversely $V_{\rm reg}>C$ suffices for a strict paid improvement, since the
state-value contribution is nonnegative.
Neither condition ranges over different probes, estimators, step schedules or
unknown covariance models.
These pre-data criteria explain both when identification can matter and when
more validation seeds cannot rescue a structurally uneconomic probe.
